\subsubsection{Minimax}

Minimax models the chasing process between the monsters and the player as a 
zero-sum game. The monsters are treated as the minimizing agent, whose goal is to 
reduce the distance to the player. The player is treated as the maximizing agent, 
whose goal is to increase the distance from the monsters.

The system builds a depth-limited game tree to represent the possible actions of 
both sides. Following the round-based structure of the game, the monsters first 
move two steps, and then the player takes one step. This process is repeated for 
a fixed number of rounds. In our implementation, the default search depth is set 
to two in order to keep the computation within the available budget.

Leaf nodes are evaluated using an evaluation function based on the distances 
between the player and all monsters:

\[
\text{Eval} = 4 \cdot \min_i d(\text{Human}, \text{Monster}_i)
+ \sum_i d(\text{Human}, \text{Monster}_i)
\]

In this function, the first term is given a weight of 4 because the nearest 
monster represents the greatest immediate risk to the player. The summation term 
encourages all monsters to continue approaching the player, rather than allowing 
only the closest monster to move effectively. A smaller evaluation score favors 
the monsters, while a larger score favors the player. If a monster reaches the 
player's cell, the node returns a large negative score. If the player is caught 
in fewer steps, the score is made even lower, encouraging the monsters to capture 
the player as quickly as possible.

Searching the entire game tree is computationally expensive, so several standard 
optimizations are applied. First, alpha-beta pruning is used to eliminate branches 
that cannot affect the final decision, thereby reducing the number of nodes that 
must be searched. Second, a transposition table is used to cache evaluated states. 
Each state is represented by the player position, the tuple of monster positions, 
and the current ply index, together with a bound flag. Since the game is played 
on a wrap-around grid, many states can be reached through different action 
sequences, so caching helps avoid repeated computation.

Third, move ordering is applied to improve pruning efficiency. Monster moves are 
ordered by increasing distance to the player, while player moves are ordered by 
decreasing distance from the monsters. In addition, the best move recorded in the 
transposition table is checked first when available. A good move ordering allows 
alpha-beta pruning to cut branches more aggressively and improves the overall 
search efficiency. After the search is completed, the selected move sequence for 
the current turn is reconstructed from the transposition table entries and 
returned through the same movement interface.

Compared with Greedy Search, Minimax provides predictive behavior. Instead of 
only following the player's current position, the monsters simulate the player's 
best possible responses and choose moves that can block escape routes. This makes 
the monsters more strategic and more difficult to avoid. However, the computational 
cost of Minimax is much higher, and it increases rapidly with the branching factor 
and search depth. Therefore, the optimizations described above are essential for 
making the algorithm practical. In this project, Minimax is used for higher 
difficulty levels where stronger monster intelligence is required.